%%%%%%%%%%%%%%%%%%%%%%%%%%%%%%%%%
% Appendix                      %
%%%%%%%%%%%%%%%%%%%%%%%%%%%%%%%%%

\setlength{\arrayrulewidth}{0.5pt}
\renewcommand{\arraystretch}{1.3}
\setlength{\tabcolsep}{4pt}
\phantomsection
\section*{Appendix A: Indicator Scoring Grids}
\addcontentsline{toc}{section}{Appendix A: Indicator Scoring Grids}

The following tables present the evidence-based scoring rubrics for all forty seven indicators of the framework, organized by dimension and sub-dimension. Each indicator is scored on a one to five scale. The CMMI level names and the interpretation of each level are as follows: Level 1 (Initial): no formal process, entirely ad hoc; Level 2 (Emerging): partial and informal practices, not repeatable; Level 3 (Defined): formalized, documented, and standardized; Level 4 (Managed): quantitative monitoring and continuous improvement; Level 5 (Optimized): systemic innovation and sector benchmarking.

\vspace{0.5em}
\noindent Indicators marked \textbf{[BCT]} are the thirteen non-skippable regulatory indicators; the flag marks a regulatory obligation whether its source is the BCT Circulaire N\textdegree{}2025-08 or the BCBS 239 principles, and Appendix~B gives the specific source of each. Indicators in D5 are \textbf{[Proxy]} indicators measured through objective organizational evidence rather than direct self-assessment.

% ─────────────────────────────────────────────────────────
% D1: Governance (25%)
% ─────────────────────────────────────────────────────────

\subsection*{A.1, D1: Governance (25\%, 3 sub-dimensions, 12 indicators)}

The Governance dimension carries the heaviest weight and the largest share of the regulatory indicators. Its twelve indicators, grouped into three sub-dimensions, are scored with the following rubrics.

\subsubsection*{Sub-dimension 1.1: Strategy and Data Policy}

\begin{center}
{\small
\begin{longtable}{|p{2.4cm}|p{2.0cm}|p{2.0cm}|p{2.2cm}|p{2.2cm}|p{2.2cm}|}
\hline
\textbf{Indicator} & \textbf{L1: Initial} & \textbf{L2: Emerging} & \textbf{L3: Defined} & \textbf{L4: Managed} & \textbf{L5: Optimized} \\
\hline
\endfirsthead
\multicolumn{6}{c}{\small (Table~\ref{tab:d1-1} continued)} \\
\hline
\textbf{Indicator} & \textbf{L1} & \textbf{L2} & \textbf{L3} & \textbf{L4} & \textbf{L5} \\
\hline
\endhead
\hline
\multicolumn{6}{r}{\small (continued)} \\
\endfoot
\hline
\caption{Dimension 1.1: Strategy and data policy scoring rubrics} \label{tab:d1-1} \\
\endlastfoot
Data strategy \textbf{[BCT]} &
  No strategy document exists &
  Informal or verbal only &
  Written, formally approved, and dated &
  Reviewed annually with KPIs &
  Integrated with corporate strategy \\
\hline
Executive sponsor &
  No designated role &
  Informal champion only &
  Formal CDO or equivalent appointed &
  CDO with dedicated budget and authority &
  CDO at C-suite with board visibility \\
\hline
Data roadmap &
  None &
  Exists but no milestones &
  Milestones defined and tracked &
  Roadmap linked to business KPIs &
  Dynamic roadmap with value metrics \\
\hline
Policy review \textbf{[BCT]} &
  Never reviewed &
  Ad hoc reviews only &
  Annual review scheduled &
  Bi-annual review with change log &
  Continuous review embedded in operations \\
\hline
\end{longtable}
}
\end{center}

\subsubsection*{Sub-dimension 1.2: Ownership and Responsibilities}

\begin{center}
{\small
\begin{longtable}{|p{2.4cm}|p{2.0cm}|p{2.0cm}|p{2.2cm}|p{2.2cm}|p{2.2cm}|}
\hline
\textbf{Indicator} & \textbf{L1: Initial} & \textbf{L2: Emerging} & \textbf{L3: Defined} & \textbf{L4: Managed} & \textbf{L5: Optimized} \\
\hline
\endfirsthead
\multicolumn{6}{c}{\small (Table~\ref{tab:d1-2} continued)} \\
\hline
\textbf{Indicator} & \textbf{L1} & \textbf{L2} & \textbf{L3} & \textbf{L4} & \textbf{L5} \\
\hline
\endhead
\hline
\multicolumn{6}{r}{\small (continued)} \\
\endfoot
\hline
\caption{Dimension 1.2: Ownership and responsibilities scoring rubrics} \label{tab:d1-2} \\
\endlastfoot
Data owners \textbf{[BCT]} &
  None identified &
  1 to 2 domains covered &
  All major domains have named owners &
  Owners trained and actively monitoring &
  Owners accountable with performance metrics \\
\hline
Data stewards &
  None &
  Informal stewardship only &
  Stewards formally designated &
  Stewards operating with defined SLAs &
  Stewards embedded in all processes \\
\hline
Accountability (RACI) &
  None documented &
  Partial documentation &
  RACI exists and communicated &
  RACI enforced and audited &
  RACI integrated into HR processes \\
\hline
Data committee &
  None &
  Informal meetings only &
  Formal committee with charter &
  Regular meetings, decisions logged &
  Committee drives cross-organization strategy \\
\hline
\end{longtable}
}
\end{center}

\subsubsection*{Sub-dimension 1.3: Regulatory Compliance (BCT and BCBS 239)}

\begin{center}
{\small
\begin{longtable}{|p{2.4cm}|p{2.0cm}|p{2.0cm}|p{2.2cm}|p{2.2cm}|p{2.2cm}|}
\hline
\textbf{Indicator} & \textbf{L1: Initial} & \textbf{L2: Emerging} & \textbf{L3: Defined} & \textbf{L4: Managed} & \textbf{L5: Optimized} \\
\hline
\endfirsthead
\multicolumn{6}{c}{\small (Table~\ref{tab:d1-3} continued)} \\
\hline
\textbf{Indicator} & \textbf{L1} & \textbf{L2} & \textbf{L3} & \textbf{L4} & \textbf{L5} \\
\hline
\endhead
\hline
\multicolumn{6}{r}{\small (continued)} \\
\endfoot
\hline
\caption{Dimension 1.3: Regulatory compliance scoring rubrics} \label{tab:d1-3} \\
\endlastfoot
BCT Circulaire mapping \textbf{[BCT]} &
  No awareness &
  Aware but not mapped &
  Requirements identified and mapped &
  Gaps identified, action plan active &
  Full compliance with evidence \\
\hline
Audit readiness \textbf{[BCT]} &
  Not ready &
  Partial documentation &
  Key documents available &
  Full audit trail maintained &
  Pre-audit self-assessment conducted \\
\hline
BCBS 239 traceability \textbf{[BCT]} &
  Not addressed &
  Partial coverage &
  Critical data elements traced &
  Full lineage documented &
  Automated lineage with controls \\
\hline
Data retention policy \textbf{[BCT]} &
  None &
  Informal rules only &
  Policy documented and approved &
  Policy enforced with monitoring &
  Policy automated and audited \\
\hline
\end{longtable}
}
\end{center}

% ─────────────────────────────────────────────────────────
% D2: Data Quality (20%)
% ─────────────────────────────────────────────────────────

\subsection*{A.2, D2: Data Quality (20\%, 3 sub-dimensions, 11 indicators)}

The Data Quality dimension is assessed through eleven indicators across three sub-dimensions, whose rubrics follow.

\subsubsection*{Sub-dimension 2.1: Quality Dimensions}

\begin{center}
{\small
\begin{longtable}{|p{2.4cm}|p{2.0cm}|p{2.0cm}|p{2.2cm}|p{2.2cm}|p{2.2cm}|}
\hline
\textbf{Indicator} & \textbf{L1: Initial} & \textbf{L2: Emerging} & \textbf{L3: Defined} & \textbf{L4: Managed} & \textbf{L5: Optimized} \\
\hline
\endfirsthead
\multicolumn{6}{c}{\small (Table~\ref{tab:d2-1} continued)} \\
\hline
\textbf{Indicator} & \textbf{L1} & \textbf{L2} & \textbf{L3} & \textbf{L4} & \textbf{L5} \\
\hline
\endhead
\hline
\multicolumn{6}{r}{\small (continued)} \\
\endfoot
\hline
\caption{Dimension 2.1: Quality dimensions scoring rubrics} \label{tab:d2-1} \\
\endlastfoot
Completeness \textbf{[BCT]} &
  Not measured &
  Ad hoc manual spot-checks &
  Measured automatically for critical datasets &
  Completeness SLAs enforced with remediation &
  Predicted via ML, proactive remediation \\
\hline
Accuracy and consistency &
  No controls &
  Manual reconciliation &
  Periodic automated reconciliation &
  Continuous automated consistency checks &
  Real-time monitoring with anomaly detection \\
\hline
Freshness and SLAs &
  No SLAs defined &
  Informal expectations &
  SLAs documented and monitored &
  SLAs reported, breach trends analyzed &
  SLAs enforced with auto-alerts and escalation \\
\hline
Deduplication &
  No deduplication &
  One-off manual deduplication &
  Defined process with matching rules &
  Automated deduplication, rate tracked &
  Real-time master data management \\
\hline
\end{longtable}
}
\end{center}

\subsubsection*{Sub-dimension 2.2: Controls and Processes}

\begin{center}
{\small
\begin{longtable}{|p{2.4cm}|p{2.0cm}|p{2.0cm}|p{2.2cm}|p{2.2cm}|p{2.2cm}|}
\hline
\textbf{Indicator} & \textbf{L1: Initial} & \textbf{L2: Emerging} & \textbf{L3: Defined} & \textbf{L4: Managed} & \textbf{L5: Optimized} \\
\hline
\endfirsthead
\multicolumn{6}{c}{\small (Table~\ref{tab:d2-2} continued)} \\
\hline
\textbf{Indicator} & \textbf{L1} & \textbf{L2} & \textbf{L3} & \textbf{L4} & \textbf{L5} \\
\hline
\endhead
\hline
\multicolumn{6}{r}{\small (continued)} \\
\endfoot
\hline
\caption{Dimension 2.2: Controls and processes scoring rubrics} \label{tab:d2-2} \\
\endlastfoot
Validation rules \textbf{[BCT]} &
  None &
  Manual checks only &
  Automated rules documented and applied in production &
  Rules version-controlled and reviewed quarterly &
  Self-learning rules, full coverage of critical flows \\
\hline
Incident tracking &
  Not tracked &
  Tracked informally &
  Formally logged and categorized &
  Tracked with root-cause analysis &
  Predictive quality monitoring \\
\hline
Remediation process &
  None &
  Ad hoc fixes only &
  Remediation process documented &
  SLA-bound remediation enforced &
  Automated remediation for known issues \\
\hline
Quality dashboards &
  None &
  Manual reports only &
  Dashboard exists and is maintained &
  Dashboard drives key decisions &
  Real-time with ML anomaly detection \\
\hline
\end{longtable}
}
\end{center}

\subsubsection*{Sub-dimension 2.3: Lineage and Traceability}

\begin{center}
{\small
\begin{longtable}{|p{2.4cm}|p{2.0cm}|p{2.0cm}|p{2.2cm}|p{2.2cm}|p{2.2cm}|}
\hline
\textbf{Indicator} & \textbf{L1: Initial} & \textbf{L2: Emerging} & \textbf{L3: Defined} & \textbf{L4: Managed} & \textbf{L5: Optimized} \\
\hline
\endfirsthead
\multicolumn{6}{c}{\small (Table~\ref{tab:d2-3} continued)} \\
\hline
\textbf{Indicator} & \textbf{L1} & \textbf{L2} & \textbf{L3} & \textbf{L4} & \textbf{L5} \\
\hline
\endhead
\hline
\multicolumn{6}{r}{\small (continued)} \\
\endfoot
\hline
\caption{Dimension 2.3: Lineage and traceability scoring rubrics} \label{tab:d2-3} \\
\endlastfoot
Data lineage \textbf{[BCT]} &
  Not documented &
  Partial and informal for some flows &
  End-to-end documented in a dedicated tool &
  Maintained automatically by ETL tooling &
  Real-time lineage, prevents reporting errors \\
\hline
Data catalog &
  None &
  Spreadsheet-based list &
  Formal catalog tool in use &
  Catalog actively maintained &
  Automated metadata harvesting \\
\hline
Metadata &
  None &
  Informal descriptions &
  Business and technical metadata defined &
  Managed in a dedicated system &
  Integrated into data governance \\
\hline
\end{longtable}
}
\end{center}

% ─────────────────────────────────────────────────────────
% D3: Architecture and Access (20%)
% ─────────────────────────────────────────────────────────

\subsection*{A.3, D3: Architecture and Access (20\%, 2 sub-dimensions, 8 indicators)}

The Architecture and Access dimension is assessed through eight indicators across two sub-dimensions, whose rubrics follow.

\subsubsection*{Sub-dimension 3.1: Centralization and Integration}

\begin{center}
{\small
\begin{longtable}{|p{2.4cm}|p{2.0cm}|p{2.0cm}|p{2.2cm}|p{2.2cm}|p{2.2cm}|}
\hline
\textbf{Indicator} & \textbf{L1: Initial} & \textbf{L2: Emerging} & \textbf{L3: Defined} & \textbf{L4: Managed} & \textbf{L5: Optimized} \\
\hline
\endfirsthead
\multicolumn{6}{c}{\small (Table~\ref{tab:d3-1} continued)} \\
\hline
\textbf{Indicator} & \textbf{L1} & \textbf{L2} & \textbf{L3} & \textbf{L4} & \textbf{L5} \\
\hline
\endhead
\hline
\multicolumn{6}{r}{\small (continued)} \\
\endfoot
\hline
\caption{Dimension 3.1: Centralization and integration scoring rubrics} \label{tab:d3-1} \\
\endlastfoot
Centralization &
  Fully siloed by department &
  Some shared datasets &
  Central repository for key data &
  Centralized with controlled exceptions &
  Unified data platform enterprise-wide \\
\hline
Interoperability &
  No integration &
  Point-to-point only &
  API standards defined &
  APIs in use across key systems &
  Full API ecosystem with governance \\
\hline
Silo reduction &
  No effort made &
  Silos identified only &
  Reduction plan in place &
  Key silos eliminated &
  Organization silo-free by design \\
\hline
Golden Record \textbf{[BCT]} &
  None &
  Identified but not built &
  Implemented for 1 to 2 critical entities &
  Implemented across all key entities &
  Enterprise-wide master data management \\
\hline
\end{longtable}
}
\end{center}

\subsubsection*{Sub-dimension 3.2: Pipelines and Infrastructure}

\begin{center}
{\small
\begin{longtable}{|p{2.4cm}|p{2.0cm}|p{2.0cm}|p{2.2cm}|p{2.2cm}|p{2.2cm}|}
\hline
\textbf{Indicator} & \textbf{L1: Initial} & \textbf{L2: Emerging} & \textbf{L3: Defined} & \textbf{L4: Managed} & \textbf{L5: Optimized} \\
\hline
\endfirsthead
\multicolumn{6}{c}{\small (Table~\ref{tab:d3-2} continued)} \\
\hline
\textbf{Indicator} & \textbf{L1} & \textbf{L2} & \textbf{L3} & \textbf{L4} & \textbf{L5} \\
\hline
\endhead
\hline
\multicolumn{6}{r}{\small (continued)} \\
\endfoot
\hline
\caption{Dimension 3.2: Pipelines and infrastructure scoring rubrics} \label{tab:d3-2} \\
\endlastfoot
Pipeline maturity &
  Manual data moves &
  Basic automated pipelines &
  Documented and monitored pipelines &
  Resilient pipelines with alerting &
  Self-healing pipelines with CI/CD \\
\hline
Access time SLA &
  No SLA, ad hoc access &
  SLA defined but not measured &
  SLA measured and reported &
  SLA consistently met &
  Real-time access for all critical data \\
\hline
Availability &
  Frequent outages &
  Basic monitoring &
  SLA-based availability targets &
  High availability with failover &
  Active-active redundancy, zero downtime \\
\hline
Access rights (RBAC) \textbf{[BCT]} &
  No controls &
  Basic user management &
  RBAC implemented with audit log &
  RBAC with full audit trails, annual review &
  Zero-trust model with anomaly detection \\
\hline
\end{longtable}
}
\end{center}

% ─────────────────────────────────────────────────────────
% D4: Analytics and Tools (20%)
% ─────────────────────────────────────────────────────────

\subsection*{A.4, D4: Analytics and Tools (20\%, 2 sub-dimensions, 8 indicators)}

The Analytics and Tools dimension is assessed through eight indicators across two sub-dimensions, whose rubrics follow.

\subsubsection*{Sub-dimension 4.1: Tool Maturity}

\begin{center}
{\small
\begin{longtable}{|p{2.4cm}|p{2.0cm}|p{2.0cm}|p{2.2cm}|p{2.2cm}|p{2.2cm}|}
\hline
\textbf{Indicator} & \textbf{L1: Initial} & \textbf{L2: Emerging} & \textbf{L3: Defined} & \textbf{L4: Managed} & \textbf{L5: Optimized} \\
\hline
\endfirsthead
\multicolumn{6}{c}{\small (Table~\ref{tab:d4-1} continued)} \\
\hline
\textbf{Indicator} & \textbf{L1} & \textbf{L2} & \textbf{L3} & \textbf{L4} & \textbf{L5} \\
\hline
\endhead
\hline
\multicolumn{6}{r}{\small (continued)} \\
\endfoot
\hline
\caption{Dimension 4.1: Tool maturity scoring rubrics} \label{tab:d4-1} \\
\endlastfoot
BI platform &
  Excel and manual tools only &
  Basic BI tool in use &
  Modern BI platform deployed &
  BI integrated enterprise-wide &
  Cloud-native platform with AI \\
\hline
Analytical coverage &
  Less than 25\% of needs covered &
  25 to 50\% &
  50 to 75\% &
  75 to 90\% &
  All analytical needs with self-service \\
\hline
Automation level &
  Fully manual &
  Some automation &
  Key processes automated &
  End-to-end for critical flows &
  ML-driven intelligent automation \\
\hline
AI and ML readiness &
  No plans &
  Awareness only &
  Pilot projects underway &
  AI/ML in production for selected use cases &
  AI/ML embedded in core processes \\
\hline
\end{longtable}
}
\end{center}

\subsubsection*{Sub-dimension 4.2: Data Usage}

\begin{center}
{\small
\begin{longtable}{|p{2.4cm}|p{2.0cm}|p{2.0cm}|p{2.2cm}|p{2.2cm}|p{2.2cm}|}
\hline
\textbf{Indicator} & \textbf{L1: Initial} & \textbf{L2: Emerging} & \textbf{L3: Defined} & \textbf{L4: Managed} & \textbf{L5: Optimized} \\
\hline
\endfirsthead
\multicolumn{6}{c}{\small (Table~\ref{tab:d4-2} continued)} \\
\hline
\textbf{Indicator} & \textbf{L1} & \textbf{L2} & \textbf{L3} & \textbf{L4} & \textbf{L5} \\
\hline
\endhead
\hline
\multicolumn{6}{r}{\small (continued)} \\
\endfoot
\hline
\caption{Dimension 4.2: Data usage scoring rubrics} \label{tab:d4-2} \\
\endlastfoot
Data-driven decisions &
  Decisions based on intuition &
  Data used occasionally &
  Data used for most major decisions &
  Data mandatory for strategic calls &
  Data-first by policy organization-wide \\
\hline
Self-service analytics &
  All reports through IT &
  Limited self-service &
  Self-service available to key teams &
  Broad self-service adoption &
  Full data democratization \\
\hline
Report utilization &
  Rarely accessed &
  Monthly usage &
  Weekly usage by key teams &
  Daily usage by most teams &
  Real-time, embedded in workflows \\
\hline
Regulatory reporting \textbf{[BCT]} &
  Fully manual &
  Partially automated, no validation trail &
  Key reports automated with validation trail &
  All reports automated, output auto-reviewed &
  End-to-end automated with real-time validation \\
\hline
\end{longtable}
}
\end{center}

% ─────────────────────────────────────────────────────────
% D5: Skills and Culture (15%)
% ─────────────────────────────────────────────────────────

\subsection*{A.5, D5: Skills and Culture (15\%, 2 sub-dimensions, 8 proxy indicators)}

The Skills and Culture dimension is assessed indirectly, through eight proxy indicators across two sub-dimensions that can be verified against objective organizational evidence. Their rubrics follow.

\noindent \textit{All indicators in D5 are proxy indicators. They are assessed exclusively through objective, auditable organizational evidence (budgets, headcount records, certification logs, and program documentation) rather than direct self-assessment.}

\subsubsection*{Sub-dimension 5.1: Data Skills (Proxy)}

\begin{center}
{\small
\begin{longtable}{|p{2.4cm}|p{2.0cm}|p{2.0cm}|p{2.2cm}|p{2.2cm}|p{2.2cm}|}
\hline
\textbf{Indicator} & \textbf{L1: Initial} & \textbf{L2: Emerging} & \textbf{L3: Defined} & \textbf{L4: Managed} & \textbf{L5: Optimized} \\
\hline
\endfirsthead
\multicolumn{6}{c}{\small (Table~\ref{tab:d5-1} continued)} \\
\hline
\textbf{Indicator} & \textbf{L1} & \textbf{L2} & \textbf{L3} & \textbf{L4} & \textbf{L5} \\
\hline
\endhead
\hline
\multicolumn{6}{r}{\small (continued)} \\
\endfoot
\hline
\caption{Dimension 5.1: Data skills proxy indicators scoring rubrics} \label{tab:d5-1} \\
\endlastfoot
Training program \textbf{[Proxy]} &
  No training programs &
  Ad hoc training only &
  Formal program with dedicated budget &
  Ongoing curriculum with L\&D integration &
  Data academy with certification paths \\
\hline
Headcount ratio \textbf{[Proxy]} &
  Less than 0.5\% of staff &
  0.5 to 1\% &
  1 to 2\% &
  2 to 5\% &
  More than 5\% \\
\hline
Certifications (past 12 months) \textbf{[Proxy]} &
  None &
  1 to 2 &
  3 to 5 &
  6 to 10 &
  More than 10 \\
\hline
Data hires (past 12 months) \textbf{[Proxy]} &
  None &
  1 to 2 hires &
  3 to 5 hires &
  6 to 10 hires &
  Dedicated data talent pipeline \\
\hline
\end{longtable}
}
\end{center}

\subsubsection*{Sub-dimension 5.2: Culture and Adoption (Proxy)}

\begin{center}
{\small
\begin{longtable}{|p{2.4cm}|p{2.0cm}|p{2.0cm}|p{2.2cm}|p{2.2cm}|p{2.2cm}|}
\hline
\textbf{Indicator} & \textbf{L1: Initial} & \textbf{L2: Emerging} & \textbf{L3: Defined} & \textbf{L4: Managed} & \textbf{L5: Optimized} \\
\hline
\endfirsthead
\multicolumn{6}{c}{\small (Table~\ref{tab:d5-2} continued)} \\
\hline
\textbf{Indicator} & \textbf{L1} & \textbf{L2} & \textbf{L3} & \textbf{L4} & \textbf{L5} \\
\hline
\endhead
\hline
\multicolumn{6}{r}{\small (continued)} \\
\endfoot
\hline
\caption{Dimension 5.2: Culture and adoption proxy indicators scoring rubrics} \label{tab:d5-2} \\
\endlastfoot
Data literacy \textbf{[Proxy]} &
  None &
  One-off awareness sessions &
  Structured program exists &
  Program reaches all staff levels &
  Literacy embedded in onboarding \\
\hline
Change management \textbf{[Proxy]} &
  No change management &
  Informal communication only &
  Change plan documented &
  Change program actively run &
  Embedded in operating model \\
\hline
Business engagement \textbf{[Proxy]} &
  Data seen as IT responsibility only &
  Occasional business input &
  Business data champions exist &
  Business co-owns data initiatives &
  Business drives the data agenda \\
\hline
Communities of practice \textbf{[Proxy]} &
  None &
  Informal groups only &
  Formal community with regular meetings &
  Active community with shared outputs &
  Multi-community ecosystem \\
\hline
\end{longtable}
}
\end{center}

\vspace{1em}
\noindent \textbf{Note on the evidence cap:} Any indicator scored above 2 out of 5 without supporting documentary evidence is automatically capped at 2 by the scoring engine, uniformly across all five dimensions. For the D5 proxy indicators the required evidence is objective and auditable, for example headcount ratios and certification counts drawn from HR records, but it must still be recorded for the score to exceed 2.

\clearpage
\phantomsection
\section*{Appendix B: Regulatory Requirements Mapping}
\addcontentsline{toc}{section}{Appendix B: Regulatory Requirements Mapping}

This appendix maps each of the thirteen non-skippable regulatory indicators to the specific obligation of the BCT Circulaire N\textdegree{}2025-08 or the BCBS~239 principles that it operationalizes. These indicators are flagged as \textbf{[BCT]} in Appendix~A and cannot be skipped in any assessment. They are concentrated in Governance (seven), Data Quality (three), Architecture and Access (two), and Analytics and Tools (one); Skills and Culture carries none.

\vspace{0.6em}
{\small
\setlength{\tabcolsep}{4pt}
\begin{longtable}{|p{3.4cm}|p{1.3cm}|p{9.6cm}|}
\hline
\textbf{Regulatory indicator} & \textbf{Dim.} & \textbf{Regulatory source and obligation} \\
\hline
\endfirsthead
\hline
\textbf{Regulatory indicator} & \textbf{Dim.} & \textbf{Regulatory source and obligation} \\
\hline
\endhead
\hline
\endfoot
\hline
\caption{Mapping of the thirteen regulatory indicators to their BCT and BCBS 239 obligations}\label{tab:reg-mapping}\\
\endlastfoot
Data strategy & D1.1 & BCT: a data strategy formally approved by executive management, dated, and made accessible; an informal or verbal strategy does not satisfy the obligation. \\
\hline
Policy review & D1.1 & BCT: a periodic, scheduled review of the data policy with a change log. \\
\hline
Data owners & D1.2 & BCT: named individual data owners with documented accountability for each critical business domain, supported by a communicated accountability matrix. \\
\hline
BCT Circulaire mapping & D1.3 & BCT: a formal, article-by-article mapping of the institution's compliance obligations to controls and evidence. \\
\hline
Audit readiness & D1.3 & BCT: a pre-assembled evidence package allowing the institution to produce supporting documents, system logs, and named responsible owners within a short delay for inspection. \\
\hline
BCBS 239 traceability & D1.3 & BCBS 239: documented and auditable lineage of every critical risk data element from its source to the regulatory report, kept current. \\
\hline
Data retention policy & D1.3 & BCT: a documented and approved policy for the classification and retention of information. \\
\hline
Completeness & D2.1 & BCBS 239: completeness of critical data across the relevant risk exposures, measured and monitored rather than assumed. \\
\hline
Validation rules & D2.2 & BCT: automated data validation rules, documented with an owner and a last review date, actively applied to data flows in production. \\
\hline
Data lineage & D2.3 & BCBS 239: end-to-end lineage and traceability for critical data flows, sufficient to reconstruct any reported figure. \\
\hline
Golden Record & D3.1 & BCT: a single source of truth, or golden record, for critical business entities, avoiding conflicting versions across systems. \\
\hline
Access rights (RBAC) & D3.1 & BCT: role-based access control with logged and periodically reviewed access grants and complete audit trails. \\
\hline
Regulatory reporting & D4.2 & BCT: automated production of regulatory reports to the BCT with a complete validation trail from source data to submitted output. \\
\hline
\end{longtable}
}

\clearpage
\phantomsection
\section*{Appendix C: Role and Permission Matrix}
\addcontentsline{toc}{section}{Appendix C: Role and Permission Matrix}

This appendix consolidates, for the delivered multi-tenant platform, which of the four roles may perform each platform capability. The roles form a strict rank order, EY Admin (owner) above Super Admin above Admin above Analyst, and every invitation creates an account exactly one rank below the inviter. The EY Admin has no bank of its own and oversees all banks; the Super Admin and Admin are scoped to a single bank; the Analyst contributes to assessments within a bank. A check mark ($\checkmark$) means the role may perform the capability; a dash (--) means it may not.

\vspace{0.6em}
{\small
\setlength{\tabcolsep}{4pt}
\begin{center}
\begin{tabular}{|p{7.2cm}|c|c|c|c|}
\hline
\textbf{Capability} & \textbf{EY Admin} & \textbf{Super Admin} & \textbf{Admin} & \textbf{Analyst} \\
\hline
Edit the EY master template & $\checkmark$ & -- & -- & -- \\
\hline
Review submissions across banks & $\checkmark$ & -- & -- & -- \\
\hline
Invite a user (exactly one rank down) & $\checkmark$ & $\checkmark$ & $\checkmark$ & -- \\
\hline
Manage users, roles, and accounts (own bank) & -- & $\checkmark$ & $\checkmark$ & -- \\
\hline
Manage the bank's departments & -- & $\checkmark$ & $\checkmark$ & -- \\
\hline
Edit the bank's questionnaire copy & -- & -- & $\checkmark$ & -- \\
\hline
Read the bank's questionnaire content & -- & $\checkmark$ & $\checkmark$ & $\checkmark$ \\
\hline
Map dimensions and launch a group assessment & -- & $\checkmark$ & $\checkmark$ & -- \\
\hline
Finalize and submit a group assessment & -- & $\checkmark$ & $\checkmark$ & -- \\
\hline
Review the bank's submissions & -- & $\checkmark$ & $\checkmark$ & -- \\
\hline
Run a solo assessment & -- & -- & -- & $\checkmark$ \\
\hline
Contribute the dimensions assigned to a department & -- & -- & -- & $\checkmark$ \\
\hline
View results, gap analysis, and compliance & -- & -- & -- & $\checkmark$ \\
\hline
Manage own account (password, language, avatar) & $\checkmark$ & $\checkmark$ & $\checkmark$ & $\checkmark$ \\
\hline
\end{tabular}
\end{center}
}
\captionof{table}{Role and permission matrix of the multi-tenant platform}
\label{tab:role-matrix}

\vspace{0.8em}
\noindent \textbf{Enforcement.} These permissions are not enforced by the interface alone. Every privileged operation is re-checked on the server, which re-reads the caller's role and bank from the database on each request, and the database itself confines every row to its own bank through row-level security, so a client that tampers with a request cannot exceed the caller's role or reach another bank's data.

\clearpage
\phantomsection
\section*{Appendix D: Platform Data Model and Tenant Isolation}
\addcontentsline{toc}{section}{Appendix D: Platform Data Model and Tenant Isolation}

This appendix summarizes the eight tables of the platform's Postgres database and the mechanism by which each bank's data is isolated from every other bank's. The tenant key is the bank name itself: there is no separate bank table, and every scoped row carries the bank name it belongs to. An empty bank name denotes the EY master template, from which each bank's questionnaire copy is seeded. Row-level security policies confine every query to rows whose bank name matches the caller's, with the EY Admin (owner) excepted by design so that it can oversee all banks.

\vspace{0.6em}
{\small
\setlength{\tabcolsep}{4pt}
\begin{center}
\begin{tabular}{|p{3.2cm}|p{8.0cm}|p{2.8cm}|}
\hline
\textbf{Table} & \textbf{Purpose} & \textbf{Tenant scope} \\
\hline
profiles & User accounts with role, bank name, department, disabled flag, and invitation lineage & Bank name \\
\hline
dimensions & Per-bank copy of the five weighted dimensions & Bank name (empty = master) \\
\hline
indicators & Per-bank copy of the forty seven indicators, with regulatory flags and rubrics & Bank name (empty = master) \\
\hline
departments & The bank's departments & Bank name \\
\hline
assessments & Group assessments, each draft or finalized, with a target level & Bank name \\
\hline
assessment\_assignments & Assignment of each dimension to exactly one department & Via the assessment \\
\hline
assessment\_answers & The shared group answers, one row per indicator, recording the contributing author & Via the assessment \\
\hline
submissions & Immutable finalized submissions, snapshotting the scores, answers, and profile & Bank name \\
\hline
\end{tabular}
\end{center}
}
\captionof{table}{The eight platform tables and their tenant scope}
\label{tab:data-model}

\vspace{0.8em}
\noindent \textbf{Isolation guarantees.} Row-level security is enabled on all eight tables. A partial unique index enforces at most one open draft assessment per bank. Submissions are immutable once written, so a finalized record cannot be altered. In-progress solo answers are the only assessment data that never leave the assessor's browser; everything else is confined to its bank at the data layer rather than by the trust of the client.

\clearpage
\phantomsection
\section*{Appendix E: Scoring Rules Quick Reference}
\addcontentsline{toc}{section}{Appendix E: Scoring Rules Quick Reference}

This appendix collects, on a single page, the exact rules by which an assessment is scored, so that any figure the tool displays can be retraced by hand.

\vspace{0.4em}
\noindent \textbf{Indicator score.} Each indicator is scored from 1 to 5 against its five-level rubric (Appendix~A).

\noindent \textbf{Evidence cap.} Any indicator scored 3 or more without supporting documentary evidence is computed as an effective score of 2 before any aggregation.

\noindent \textbf{Skip logic.} At most $\lfloor 0.20 \times n \rfloor$ indicators of a dimension of $n$ indicators may be skipped, giving ceilings of 2 (Governance), 2 (Data Quality), 1 (Architecture and Access), 1 (Analytics and Tools), and 1 (Skills and Culture). Skipped indicators are excluded from the denominator, not scored as zero. The thirteen regulatory indicators can never be skipped.

\noindent \textbf{Aggregation.} A sub-dimension score is the average of the effective scores of its indicators; a dimension score is the average of its sub-dimension scores; the Global Maturity Index (GMI) is the weighted sum of the dimension scores, with weights 25/20/20/20/15, renormalized over the weights of the answered dimensions only.

\noindent \textbf{Percentage.} The percentage form is $\mathrm{round}(\mathrm{GMI} \times 20)$.

\vspace{0.5em}
\noindent \textbf{Maturity bands.} \quad
Level 1 Initial: 1.00 to 1.79 \quad$\vert$\quad
Level 2 Emerging: 1.80 to 2.59 \quad$\vert$\quad
Level 3 Defined: 2.60 to 3.39 \quad$\vert$\quad
Level 4 Managed: 3.40 to 4.19 \quad$\vert$\quad
Level 5 Optimized: 4.20 to 5.00.

\vspace{0.5em}
\noindent \textbf{Regulatory compliance.} A regulatory indicator is compliant when its effective score reaches 3, the point at which an obligation is formalized, documented, and demonstrable rather than merely intended, non-compliant below 3, and pending when unanswered. The compliance rate is the share of the thirteen regulatory indicators that are compliant. The regulatory exposure is Low when the rate reaches 80\%, Medium when it reaches 50\%, and High otherwise; these exposure bands are an operational risk signal chosen for the tool, not a threshold defined by the regulator, and a formal compliance exercise would treat any non-compliant regulatory indicator as requiring remediation regardless of the aggregate rate.
